% =====================================================================
%  A Dual-Tank Model for Real-Time Tracking of Phosphocreatine and
%  Glycolytic Energy Reserves in Cycling
%
%  Standalone journal-style manuscript.
%
%  Compile locally (no shell-escape needed):
%     pdflatex dual-tank-anaerobic-model-journal.tex
%     pdflatex dual-tank-anaerobic-model-journal.tex      % 2nd pass: refs/xrefs
%     pdflatex dual-tank-anaerobic-model-journal.tex      % 3rd pass: cleveref
%  (three passes resolve all \ref / \cref / \cite numbering.)
%
%  Portable: uses only standard TeX Live packages, pdflatex-compatible,
%  no external fonts. For a Times-like look, uncomment the newtx lines
%  in the preamble (requires the newtx package).
% =====================================================================
\documentclass[11pt,a4paper]{article}

\usepackage[T1]{fontenc}
\usepackage[utf8]{inputenc}
\usepackage{lmodern}
% ---- optional Times-like fonts (uncomment if `newtx` is installed) ----
% \usepackage{newtxtext,newtxmath}
\usepackage{microtype}

\usepackage[a4paper,margin=1in,headheight=14pt]{geometry}
\usepackage{amsmath,amssymb}
\usepackage{siunitx}
\usepackage{xcolor}
\usepackage{booktabs}
\usepackage{array}
\usepackage{tabularx}
\usepackage{enumitem}
\usepackage{caption}
\usepackage{listings}
\usepackage{titlesec}
\usepackage{fancyhdr}

% ---- palette (project convention: purple = PCr, green = glycolytic) ----
% Defined BEFORE hyperref so allcolors=linkblue resolves.
\definecolor{pcrpurple}{HTML}{6C26A8}
\definecolor{glygreen}{HTML}{1A8C3A}
\definecolor{linkblue}{HTML}{1A5FB4}
\definecolor{rulegray}{HTML}{BBBBBB}
\definecolor{boxbg}{HTML}{F5F1FA}
\definecolor{codebg}{HTML}{F7F5FB}

\usepackage[numbers,sort&compress,super]{natbib}
\usepackage[colorlinks=true,allcolors=linkblue]{hyperref}
\usepackage[nameinlink,capitalise]{cleveref}

% ---- modern section styling ----
\titleformat{\section}{\color{pcrpurple}\Large\bfseries}{\thesection}{0.6em}{}
\titleformat{\subsection}{\color{glygreen!85!black}\large\bfseries}{\thesubsection}{0.6em}{}
\titleformat{\subsubsection}{\normalsize\bfseries\itshape}{\thesubsubsection}{0.6em}{}
\titlespacing*{\section}{0pt}{1.4em}{0.5em}

\pagestyle{fancy}\fancyhf{}
\renewcommand{\headrulewidth}{0.4pt}
\fancyhead[L]{\small\itshape Dual-Tank Anaerobic Reserve Model}
\fancyhead[R]{\small\thepage}

% ---- listings: language-agnostic pseudocode (ASCII, pdflatex-safe) ----
\lstdefinestyle{pseudo}{
  basicstyle=\ttfamily\small,
  backgroundcolor=\color{codebg},
  frame=leftline, framerule=1.4pt, rulecolor=\color{pcrpurple},
  xleftmargin=1.1em, framexleftmargin=1.1em,
  breaklines=true, columns=fullflexible, keepspaces=true,
  commentstyle=\color{glygreen!60!black}\itshape,
  morecomment=[l]{\#}, showstringspaces=false, aboveskip=0.8em, belowskip=0.8em}

\captionsetup{font=small,labelfont={bf,color=pcrpurple},labelsep=period}

% ---- key-points box ----
\newcommand{\keybox}[1]{%
  \begin{center}\fcolorbox{pcrpurple}{boxbg}{%
  \begin{minipage}{0.97\linewidth}\small #1\end{minipage}}\end{center}}

% ---- symbol macros (consistent typography throughout) ----
\newcommand{\Wp}{\ensuremath{W'}}
\newcommand{\Wbal}{\ensuremath{W'\!\text{bal}}}
\newcommand{\CP}{\ensuremath{\mathrm{CP}}}
\newcommand{\LTone}{\ensuremath{\mathrm{LT1}}}
\newcommand{\fp}{\ensuremath{f_{\mathrm p}}}
\newcommand{\Rp}{\ensuremath{R_{\mathrm p}}}
\newcommand{\Rg}{\ensuremath{R_{\mathrm g}}}
\newcommand{\Cp}{\ensuremath{C_{\mathrm p}}}
\newcommand{\Cg}{\ensuremath{C_{\mathrm g}}}
\newcommand{\Ppmax}{\ensuremath{P_{\mathrm{p,max}}}}
\newcommand{\grate}{\ensuremath{g_{\mathrm{rate}}}}
\newcommand{\gpmax}{\ensuremath{g_{\mathrm{p,max}}}}
\newcommand{\taup}{\ensuremath{\tau_{\mathrm p}}}
\newcommand{\taug}{\ensuremath{\tau_{\mathrm g}}}
\newcommand{\tauon}{\ensuremath{\tau_{\mathrm{on}}}}
\newcommand{\tauoff}{\ensuremath{\tau_{\mathrm{off}}}}
\newcommand{\tauW}{\ensuremath{\tau_{W'}}}
\newcommand{\tauPCr}{\ensuremath{\tau_{\mathrm{PCr}}}}
\newcommand{\taudep}{\ensuremath{\tau_{\mathrm{dep}}}}
\newcommand{\tauaer}{\ensuremath{\tau_{\mathrm{aer}}}}
\newcommand{\gatep}{\ensuremath{\mathrm{gate}_{\mathrm p}}}
\newcommand{\Pos}{\ensuremath{P_{1\mathrm s}}}
\newcommand{\halft}{\ensuremath{t_{1/2}}}
\newcommand{\degC}{\textsuperscript{31}P-MRS}

\hypersetup{
  pdftitle={A Dual-Tank Model for Real-Time Tracking of Phosphocreatine and Glycolytic Energy Reserves in Cycling},
  pdfauthor={S. J. Cieply},
  pdfsubject={Critical Power, W-prime, anaerobic energy systems, cycling}}

\setlength{\parskip}{0.35em}
\setlength{\parindent}{0pt}

% =====================================================================
\begin{document}

\begin{center}
  {\LARGE\bfseries A Dual-Tank Model for Real-Time Tracking of\\[2pt]
   Phosphocreatine and Glycolytic Energy Reserves in Cycling\par}
  \vspace{1.0em}
  {\large S.\,J.\ Cieply\par}
  \vspace{0.3em}
  {\normalsize AnaerobicFuelTanks Project\par}
  {\small\href{mailto:sjcieply@gmail.com}{sjcieply@gmail.com}\par}
\end{center}

\vspace{0.6em}

\begin{abstract}\noindent
Endurance-cycling analytics have converged on the Critical Power (\CP) / \Wp{} framework, in which all
work performed above a sustainable power asymptote is drawn from a single finite reserve, \Wp. This lumps
two physiologically distinct anaerobic systems---the fast phosphocreatine (PCr / alactic) system and the
slower glycolytic (lactic) system---into one tank with a single recovery time constant. That
simplification discards the information a rider most wants during hard, variable-intensity efforts:
\emph{which} reserve is being spent, and \emph{how fast} each will return. This paper develops a
\textbf{dual-tank model} that splits \Wp{} into a fast reserve and a slow reserve, drives both from the
instantaneous power trace, and applies system-specific depletion and restoration laws drawn from the
bioenergetic and \degC{} literature. The model is explicitly a physiologically-motivated decomposition of
one measured quantity (\Wbal) whose split is assumed rather than measured, and whose extra resolution over
single-tank \Wbal{} is concentrated in the transients after hard efforts. It generalises \Wbal{} in
depletion and diverges from it, by design, in recovery. The model is reduced to be computable at
1~Hz on a wrist- or bar-mounted device, and a Garmin Connect~IQ implementation displays
live reserve, consumption, and depletion for both systems. We give a validation strategy---including the
out-of-sample test that would earn the second tank and a compartment-level test the current model
fails---and an explicit statement of assumptions and limits.
\end{abstract}

\keybox{%
\textbf{Key points.}
\begin{itemize}[leftmargin=1.2em,itemsep=1pt,topsep=2pt]
\item The two bars are \textbf{compartments of \Wp}, named for their dominant metabolic system; the
metabolic identification is motivation, not a measured claim---\Wp{} recovers ${\sim}4\times$ faster than
the muscle metabolites the tanks are named after (\cref{subsec:recon}).
\item Non-redundant content lives in \emph{transients}: the glycolytic activation ramp at effort onset and
the fast-reserve recovery transient after hard efforts. During steady supra-\CP{} effort both bars are
affine in \Wbal.
\item Under an oxidative-headroom recovery gate, the effective fast-reserve recovery constant is
100--500~s while a rider is pedalling, so the fast-reserve bar is
\emph{live} through the recovery valleys of an interval session.
\item For training prescription, the cumulative per-system load separates an alactic from a glycolytic
session; the discriminating mechanism is the \textbf{glycolytic flux ceiling scaled by the activation
ramp}, the metric is blind to \emph{submaximal} alactic work, and it must be judged against a
zero-parameter effort-duration baseline, not only \Wbal{} (\cref{subsec:trainingload}).
\item The compartment-level partition is testable in-modality by muscle biopsy during cycle ergometry;
against repeated-sprint data the current model fails in level and direction, so the per-system split is a
descriptive \Wp-statistic, not a validated bioenergetic ATP partition (\cref{subsec:repeatedsprint}).
\end{itemize}}

\textbf{Keywords:} critical power; \Wp; phosphocreatine; anaerobic glycolysis; \Wbal; repeated-sprint;
real-time cycling analytics; Connect~IQ.

% =====================================================================
\section{Introduction}\label{sec:intro}

A power meter reports mechanical output, not metabolic cost. The bridge between the two is a model of the
body's energy-supply systems. The dominant bridge in cycling---the \CP/\Wp{} model---treats supra-\CP{}
work as depletion of a single reservoir and sub-\CP{} work as its exponential refill.\citep{Skiba2012} It
is power-native, validated, and computationally trivial, which is why it underlies \Wbal{} displays in
GoldenCheetah, intervals.icu, and several head units.

Its central simplification is also its central weakness for real-time decision-making: \Wp{} is one tank.
Physiologically, anaerobic capacity is at least two tanks with very different
dynamics:\citep{Margaria1933,diPrampero1999}

\begin{itemize}[leftmargin=1.4em,itemsep=1pt]
\item The \textbf{phosphocreatine (PCr / alactic)} system---high power, low capacity, \emph{fast}
oxidative recovery (seconds to a minute), with recovery that is pH-sensitive and slows across repeated
bouts.
\item The \textbf{glycolytic (lactic)} system---larger capacity, lower peak power, and \emph{slow}
recovery (minutes) that effectively proceeds only at low intensity.
\end{itemize}

A rider deciding whether to cover the next surge cares which tank is empty. If the fast reserve is full
but the slow reserve is depleted, a 5~s jump is fine but a 40~s dig
is not---a distinction single-tank \Wbal{} cannot express. The goal is a model that keeps the power-native,
on-device virtues of \Wbal{} while restoring the two-system resolution.

\paragraph{Two use cases, two outputs.} The first---\emph{pacing/racing}---uses the live bars in the
moment (``can I cover this attack?''); its value is concentrated in the post-effort transients
(\cref{sec:assumptions}) and is real but narrow. The second---\textbf{training prescription and load
monitoring}---uses the \emph{cumulative per-system load} over a session, and is where the decomposition is
most clearly worth more than \Wbal{} (\cref{subsec:trainingload}). The two anaerobic systems carry very
different physiological and recovery costs: alactic surges---short, high-power, with adequate
recovery---resynthesise in seconds to a minute and can be repeated at high volume with little systemic
fatigue; glycolytic surges accumulate H\textsuperscript{+}/P\textsubscript{i} and clear over many minutes
(lactate \halft{} $\approx$ 1366~s\citep{Ferguson2010}) and cost far more recovery.
An athlete targeting the PCr system---alactic power, repeat-sprint ability---must keep efforts short and
recoveries full enough that glycolysis stays disengaged, and needs to know, per session, how much of the
anaerobic work was alactic versus glycolytic. \Wbal{} collapses that to one number; the two-tank model
does not.

\paragraph{Two honest qualifications.} First, the split is a modelling assumption, not a measurement:
power alone cannot identify how \Wp{} divides between the two systems (\cref{subsec:update}). Second, the
extra resolution is concentrated in transients. Whenever the fast reserve is full, the slow reserve is
merely an affine rescaling of single \Wbal. The fast reserve is \emph{not} full most of the time while
riding, however: under the oxidative-headroom recovery gate (\cref{subsec:update}) its effective recovery
constant is \SIrange{100}{500}{\second} at typical soft-pedal powers rather than the \SI{27}{\second} of a
stopped rider, so the bar stays live through the recovery valleys of an interval session
(\cref{subsec:usefulness}). The model earns its keep as a heuristic decomposition of \Wp, not as two
independently measured reserves.

% =====================================================================
\section{Physiological background}\label{sec:physio}

\paragraph{Phosphocreatine (alactic).} ATP is buffered by the creatine-kinase reaction
(PCr\,+\,ADP\,$\rightarrow$\,ATP\,+\,Cr). Intramuscular PCr ($\sim$20--25~mmol\,kg\textsuperscript{-1} wet
muscle) is the highest-power energy source and is drawn on instantly at any power transient. Its
resynthesis is oxidative and fast: the classic alactic O\textsubscript{2}-debt half-time is
$\approx$\SI{30}{\second},\citep{Margaria1933} and \degC{} gives recovery time constants
$\tauPCr \approx$ \SIrange{20}{40}{\second}, strongly \textbf{pH-dependent} (acidosis slows it) and
biphasic under high glycolytic load.\citep{Harris1976,Yoshida2013} Recovery also slows progressively
across repeated hard bouts.

\paragraph{Glycolytic (lactic).} Anaerobic glycolysis converts glycogen to lactate, supplying large ATP
flux for ${\sim}10\,\mathrm{s}$--$2\,\mathrm{min}$. Its ``fuel level'' is best tracked as accumulated
muscle/blood lactate, a proxy for H\textsuperscript{+} and metabolite accumulation. Critically, glycolytic
flux itself \emph{falls} across repeated maximal sprints: acidosis inhibits phosphorylase and
phosphofructokinase, so the glycolytic ATP contribution decays from $\approx$40\% of total ATP in the
first of $10\times\SI{6}{\second}$ efforts to $<$10\% by the last,\citep{Gaitanos1993,SciRep2024} while
the aerobic contribution rises to absorb the shortfall ($\sim$29\%$\rightarrow$43\% between
sprints).\citep{Bogdanis1996} Restoration of the glycolytic reserve is slow (lactic O\textsubscript{2}-debt
half-time $\approx$~15~min) and proceeds fastest at low intensity, effectively ceasing above the first
lactate threshold (\LTone). The model as formulated has no term for glycolytic flux falling with fatigue;
this is examined directly, with an optional corrective term, in \cref{subsec:repeatedsprint}.

\paragraph{The onset that matters.} Both systems contribute from the start of a hard effort, but not
equally at first: PCr is drawn on instantly, while glycolytic flux takes several seconds to ramp up as
glycogen phosphorylase is activated,\citep{Parolin1999} so the early demand is PCr-weighted and shifts
toward a shared draw as glycolysis engages.\citep{GonzalezAlonso2000} This graded onset---and the very
different refill rates---is what the dual-tank model must reproduce (via a parallel, PCr-weighted draw and
a glycolytic activation ramp; \cref{subsec:update}), rather than a strict ``PCr first, then glycolytic''
hand-off.

\paragraph{What the two ``tanks'' represent.} The tank is a deliberate simplification for intuition and
on-device display, not a claim that either system is a literal reservoir that drains. The two sides are
not even the same \emph{kind} of quantity. The fast reserve corresponds to a real, depletable substrate
store (intramuscular phosphocreatine), so ``empty'' is close to literal. The slow reserve is better
understood as \emph{tolerance to accumulating fatigue-related metabolites} than as a fuel gauge: on the
1--4~min timescale it represents, muscle glycogen is nowhere near exhausted (carbohydrate depletion is a
separate, hours-long limiter this model does not address). What forces power down at the limit of a hard
effort is the buildup of inorganic phosphate (P\textsubscript{i}), H\textsuperscript{+}, ADP, and
extracellular K\textsuperscript{+} toward the limit of tolerance, consistent with exhaustion at \CP/\Wp{}
coinciding with a reproducible low-PCr / high-metabolite state (lactate itself is largely a fuel and a
marker, not the cause of force loss). This is why the two recovery laws differ: the fast reserve
resynthesises quickly, whereas the slow side recovers slowly and only at low intensity because it reflects
clearance and buffering of accumulated byproducts, which need aerobic metabolism and time. In short:
\textbf{the fast reserve is a fuel that depletes; the slow reserve is a byproduct bucket that fills.}

% =====================================================================
\section{Prior models}\label{sec:prior}

Three model families exist:

\begin{itemize}[leftmargin=1.4em,itemsep=2pt]
\item \textbf{Family A --- \CP/\Wp-balance.}\citep{Skiba2012,Skiba2015,Bartram2018} Power-native and
on-device-friendly, but one lumped tank with a single recovery~$\tau$. The recovery law
$\tauW = 546\,e^{-0.01\,D_{\CP}} + 316~\mathrm{s}$ cannot represent fast PCr and slow glycolytic recovery
simultaneously.
\item \textbf{Family B --- hydraulic tank models.}\citep{Morton1986,Weigend2021} Two explicit anaerobic
vessels (phosphagen, glycolytic) plus an aerobic vessel; out-predicts \Wbal{} for intermittent recovery.
Its eight parameters are abstract and fit per athlete by evolutionary computation---heavy for an embedded
target.
\item \textbf{Family C --- bioenergetic supply/demand ODEs.}\citep{Chorzempa2023} Explicit alactic /
lactic / aerobic metabolic-rate terms, with an efficiency-limited PCr recovery and a power-gated
lactate-removal law. Mechanistically ideal but has 17 parameters and needs grey-box fitting.
\end{itemize}

\textbf{The gap.} Family~A is implementable but under-resolved; Families~B and~C resolve the two systems
but are too heavy to run and calibrate on a head unit. The model below keeps the two-tank resolution of
B/C with the on-device simplicity of A.

To be precise about ``lighter'': the model's core count
($\CP, \Wp, \fp, \Ppmax, \grate, \taup, \taug, \tauon, \LTone$, with $\tauoff = \tauon$ by default---nine
free numeric parameters, plus three optional realism terms and two functional-form gates,
$\gatep$ and the intensity-dependent glycolytic recovery shape) is essentially the hydraulic model's
eight: the same size, not a smaller model. The difference is \emph{where the parameters come from}. Only
\CP{} and \Wp{} are fitted per athlete (from a standard \CP{} test); the rest are hard-coded to literature
defaults rather than fit by evolutionary computation. That buys on-device simplicity and a trivial
calibration, at a cost: the load-bearing defaults ($\fp$, $\Ppmax$, $\LTone$, and the recovery $\tau$'s)
are exactly the quantities a \CP{} test does \emph{not} determine, so the two-system resolution rests on
assumptions, not measurements. This is a legitimate identifiability-for-convenience trade, and the
sections below are explicit about which numbers are earned and which are assumed.

% =====================================================================
\section{The dual-tank model}\label{sec:model}

\subsection{State and parameters}\label{subsec:state}

Two energy reserves (Joules), each with a fixed capacity: the \textbf{fast reserve} \Rp{} (capacity \Cp)
and the \textbf{slow reserve} \Rg{} (capacity \Cg), with total anaerobic capacity $\Wp = \Cp + \Cg$. We
call them the fast reserve (``PCr'') and slow reserve (``glycolytic''); the metabolic labels are
motivation, not a measured claim. As \cref{subsec:recon} shows, the reserves are compartments of \Wp, with
the correct internal fast/slow ratio ($\sim$18$\times$) but an absolute timescale $\sim$4$\times$ faster
than the muscle metabolites they are named after, because \Wp{} itself recovers $\sim$4$\times$ faster
than PCr or lactate. The bar is therefore a fast-recovering component of \Wbal, and the model is right
about \Wp's internal dynamics while remaining agnostic about the metabolites.

\Cref{tab:params} lists the parameters. Derived quantities: $\Cp = \fp\,\Wp$, $\Cg = (1-\fp)\,\Wp$, with
initial state $\Rp = \Cp$, $\Rg = \Cg$ (full). Nine parameters are free numeric values; two
functional-form gates ($\gatep$ and the glycolytic recovery shape) and three optional realism terms
(pH-slowed PCr recovery, aerobic ramp, glycolytic flux fatigue) sit on top. Only \CP{} and \Wp{} are
fitted per athlete. A fast-reserve depletion kinetic, $\taudep = \Cp/\Ppmax$, also follows from the table
(\cref{subsec:behaviour}) and is equally assumed.

\begin{table}[htbp]\centering\small
\caption{Model parameters. Only \CP{} and \Wp{} are fitted per athlete; the rest are literature-set.}
\label{tab:params}
\renewcommand{\arraystretch}{1.25}
\begin{tabularx}{\linewidth}{@{}l>{\raggedright\arraybackslash}X>{\raggedright\arraybackslash}X@{}}
\toprule
\textbf{Symbol} & \textbf{Meaning} & \textbf{Typical default} \\
\midrule
\CP & critical power (W) & from a 3-/12-min \CP{} test ($\approx$ FTP $+$ a few \%) \\
\Wp & total work above \CP{} (J) & from the same test ($\sim$15--25~kJ) \\
\fp & fraction of \Wp{} in the fast reserve & \textbf{0.25} --- assumed (physiology; \cref{sec:assumptions}) \\
\Ppmax & fast-reserve peak power above \CP{} (W), at a \emph{full} tank & $\approx \Pos - \CP$ (best 1~s power $-$ \CP) \\
\grate & glycolytic peak flux as a fraction of fast-reserve peak flux & \textbf{0.5 --- load-bearing};\citep{diPrampero1999} band $\sim$0.3--0.7, drives the training-load separation (\cref{subsec:trainingload}) \\
\taup & fast-reserve recovery time constant (s); a \Wp-recovery constant, not muscle-PCr resynthesis & \textbf{27} --- set by the \Wp-reconstitution curve (\cref{subsec:recon}); do not read as biopsy PCr \\
\taug & slow-reserve recovery time constant (s) & \textbf{470} --- reconstitution fit\citep{Ferguson2010,ChorleyLamb2020} (\cref{subsec:recon}) \\
\tauon & glycolytic activation time constant (s) & 6\citep{Parolin1999} \\
\tauoff & glycolytic de-activation time constant (s) & $=\tauon$ --- assumed; not load-bearing (\cref{subsec:trainingload}) \\
\LTone & first lactate threshold power (W); sets the band of the recovery shape & \emph{measured} (a threshold test), not a fixed \%\CP \\
\gatep & oxidative-headroom recovery gate $(\CP-P)/\CP$ --- a functional form, no free constant & assumed; shape is a choice (\cref{sec:assumptions}) \\
$g_{\mathrm{fat}}$ & \emph{optional} glycolytic flux-fatigue exponent, $\mathrm{rate}_g\,(\Rg/\Cg)^{g_{\mathrm{fat}}}$ & \textbf{0 (off)} by default (\cref{subsec:repeatedsprint}) \\
\bottomrule
\end{tabularx}
\end{table}

\paragraph{On \fp.} We set \textbf{0.25} from physiological alactic-fraction estimates (the PCr system
supplies $\sim$20--30\% of anaerobic ATP capacity\citep{diPrampero1999,Bangsbo1990}); once the
reconstitution reference protocol's recovery power is known (20~W, near-passive), the reconstitution curve
independently implies $\fp \approx 0.20$--$0.25$. Treat it as a per-athlete calibration target on the band
$\sim$0.20--0.25.

\paragraph{On \Ppmax.} $\Ppmax \approx \Pos - \CP$ is the rate at a \emph{full} tank. Fast-reserve flux is
not constant: the rate \emph{ceiling} tapers with tank fullness (creatine-kinase equilibrium---flux falls
as the store depletes), so the available rate is $\Ppmax\,(\Rp/\Cp)$, equal to \Ppmax{} only when the tank
is full (\cref{subsec:update}). Read $\Pos - \CP$ as an upper bound: at 1~s glycolysis is already $\sim$15\%
active, so it mildly over-attributes to the fast reserve.

\paragraph{On \grate.} The ratio $\gpmax = 0.5\,\Ppmax$ is not arbitrary: peak glycolytic ATP flux is
roughly half peak PCr/creatine-kinase flux in human muscle,\citep{diPrampero1999,Parolin1999} so the fast
system is higher-power by about 2:1. \grate{} is nonetheless load-bearing: the ablation in
\cref{subsec:trainingload} shows it (via the glycolytic flux ceiling) generates the training-load
separation, and sweeping it 0.3$\rightarrow$0.7 moves that separation from $\sim$48 to $\sim$33 points. It
is banded ($\sim$0.3--0.7) and should be checked against sprint biopsy data (\cref{subsec:repeatedsprint}),
against which it has never been fit.

\subsection{Reconstitution: calibration source and a tension it raises}\label{subsec:recon}

The 37\,/\,65\,/\,86\% \Wp-recovery at 2\,/\,6\,/\,15~min (half-time \SI{234}{\second}) is from
Ferguson et~al.,\citep{Ferguson2010} reviewed by Chorley \& Lamb,\citep{ChorleyLamb2020} who state it was
measured when recovering at a nominal 20~W. Three things follow.

\textbf{(1) \taug{} is calibrated to $\sim$470~s.} The recovery gate makes the effective constant depend on
recovery power; the 20~W figure resolves this. But 20~W is not quite $\mathrm{gate}=1$: with
$\LTone \approx 0.8\,\CP \approx 200~\mathrm{W}$, $\mathrm{gate} = (200-20)/200 = 0.90$, so a model
reproducing Ferguson needs $\taug \approx 470~\mathrm{s}$.

\textbf{(2) \fp{} is consistent with the curve, not independently measured by it.} With recovery power
fixed, the curve constrains a ridge in $(\fp,\taug)$, not \fp{} alone: both $(\fp\,0.25, \taug\,520)$ and
the joint optimum $(\fp\,0.20, \taug\,470)$ fit 37/65/86 to within $\sim$3--5 points. Physiology selects
$\fp \approx 0.25$ on that ridge; the recovery data alone would prefer $\sim$0.20.

\textbf{(3) The tanks are compartments of \Wp, not of muscle metabolites.} Ferguson reports the two
recovery channels underlying \Wp{} reconstitution: VO\textsubscript{2} $\halft = \SI{74}{\second}$ (a proxy
for oxidative/PCr resynthesis) and blood-lactate $\halft = \SI{1366}{\second}$. Placing these beside the
model's component half-times (\cref{tab:components}) shows the model's two constants are a near-uniform
$\sim$4$\times$ rescaling of Ferguson's channels---the \emph{shape} (the 18-fold fast/slow separation) is
reproduced almost exactly, only the \emph{scale} differs, and by the same factor in both. That is not the
signature of fitting two exponentials to a sum. The offset is unavoidable: building the two-compartment
model from Ferguson's \emph{measured} channel constants (107~s, 1971~s) cannot reproduce Ferguson's own
recovery curve (\cref{tab:fergusonbuild})---a metabolite-calibrated model predicts half the observed
recovery. \Wp{} simply recovers $\sim$4$\times$ faster than the PCr/lactate it correlates with, consistent
with Ferguson's conclusion that \Wp{} reconstitution is ``not a unique function of phosphocreatine
concentration or arterial [lactate]\dots unlikely to simply reflect a finite energy store.'' What differs
from the metabolites is therefore the \emph{naming}, not the time constants. Symmetrically, blood lactate
\halft{} is no more a measure of \emph{muscle} H\textsuperscript{+}/P\textsubscript{i} clearance than
VO\textsubscript{2} off-kinetics is a measure of PCr resynthesis, so neither channel measures a tank and
neither adjudicates the decomposition. The fast/slow-reserve naming is the resolution.

\begin{table}[htbp]\centering\small
\caption{The model's recovery components against Ferguson's measured channels: the ratio is preserved, the
absolute scale is $\sim$4$\times$ faster.}
\label{tab:components}
\renewcommand{\arraystretch}{1.2}
\begin{tabular}{@{}llll@{}}
\toprule
\textbf{Component} & \textbf{Model} & \textbf{Ferguson channel} & \textbf{Ratio} \\
\midrule
fast & $\taup = 27~\mathrm{s} \rightarrow \halft = 18.7~\mathrm{s}$ & VO\textsubscript{2} 74~s & 4.0$\times$ \\
slow & $\taug = 470~\mathrm{s} \rightarrow \halft = 326~\mathrm{s}$ & lactate 1366~s & 4.2$\times$ \\
\textbf{ratio (slow/fast)} & \textbf{17.4$\times$} & \textbf{18.5$\times$} & within 6\% \\
\bottomrule
\end{tabular}
\end{table}

\begin{table}[htbp]\centering\small
\caption{A model built from Ferguson's own channel constants under-predicts Ferguson's own recovery curve
by roughly half; the calibrated model reproduces it.}
\label{tab:fergusonbuild}
\renewcommand{\arraystretch}{1.2}
\begin{tabular}{@{}lccc@{}}
\toprule
 & \textbf{2~min} & \textbf{6~min} & \textbf{15~min} \\
\midrule
Observed\citep{Ferguson2010} & 37\% & 65\% & 86\% \\
Model ($\taup\,27$, $\taug\,470$) & 40\% & 63\% & 87\% \\
Built from Ferguson's channels (107~s, 1971~s) & 20\% & 35\% & 50\% \\
\bottomrule
\end{tabular}
\end{table}

This reconstitution tension does not reach the training-load result of \cref{subsec:trainingload}: that
signal is driven by the \emph{depletion} law, anchored in measured activation kinetics (\tauon{} from
phosphorylase activation\citep{Parolin1999}) and the flux ratio (\grate\citep{diPrampero1999}), not by the
recovery constants. Substituting Ferguson's own $\taup = \SI{107}{\second}$ moves the alactic/glycolytic
separation by $<$1 point.

\subsection{The 1~Hz update}\label{subsec:update}

Let $P$ be instantaneous power (W) and $\Delta t = 1~\mathrm{s}$; define demand relative to aerobic supply
as $\Delta = P - \CP$.

\paragraph{Case 1 --- depletion ($\Delta > 0$, above \CP).} Demand $\mathrm{need} = \Delta\,\Delta t$ (J)
is met by both systems in parallel. Glycolytic supply is not instantaneous: glycogen phosphorylase
transforms to its active form over the first several seconds of maximal effort (rising from 12\% at rest to
47\% by 6~s\citep{Parolin1999}), while phosphocreatine is the immediate buffer.\citep{Bogdanis1996,Bogdanis1998}
Both contribute from onset, not in sequence.\citep{GonzalezAlonso2000} A glycolytic activation term
$g \in [0,1]$ ramps first-order with time constant $\tauon \approx 6~\mathrm{s}$ above \CP{} and relaxes
during recovery. The share of submaximal demand is apportioned by \emph{capacity}; a separate peak-flux
\emph{ceiling}, with the fast-reserve rate tapered by fullness, governs maximal efforts:

\begin{lstlisting}[style=pseudo]
g       <- g + (1 - g)*(1 - exp(-dt/tau_on))    # glycolytic activation ("inertia")

# --- SHARE of submaximal demand: capacity-proportional (not rate-proportional) ---
w_p     = C_p                                   # g=0 -> fast covers all; g=1 -> fast covers f_p
w_g     = C_g * g
share_p = need * w_p / (w_p + w_g)              # guard: if w_p+w_g ~ 0, share_p = need
share_g = need - share_p

# --- RATE CEILING: peak flux, fast-reserve tapered with fullness (maximal efforts only) ---
rate_p  = P_p_max * (R_p / C_p)                 # guard: C_p > 0
rate_g  = g_rate  * P_p_max * g

take_p  = min(share_p, R_p, rate_p*dt)          # each reserve capacity- AND rate-limited
take_g  = min(share_g, R_g, rate_g*dt)
unmet   = need - take_p - take_g
take_g += min(unmet, R_g - take_g, rate_g*dt - take_g);  unmet -= ...   # spill to partner,
take_p += min(unmet, R_p - take_p, rate_p*dt - take_p);  unmet -= ...   # then to a deficit
R_p    -= take_p;  R_g -= take_g
D      += unmet                                 # DEFICIT: supra-CP work the caps could not place
if unmet > 0:  rate_limited = true              # power beyond the reserves' flux (often a stale P_1s)
\end{lstlisting}

\emph{Why the share and the ceiling are different objects.} The ceiling ($\Ppmax$, $\gpmax$) encodes ``the
fast system is higher-power,'' true at maximal effort where both systems are flux-limited. It is the wrong
object for apportioning \emph{submaximal} supra-\CP{} demand, where neither system is near its flux limit
and sharing is set by metabolic control. Weighting the share by capacity while keeping the taper on the
ceiling gives the correct behaviour in both regimes:

\begin{itemize}[leftmargin=1.4em,itemsep=2pt]
\item \textbf{Submaximal:} both reserves drain $\approx$proportionally to capacity, so
$\Rp/\Cp \approx \Rg/\Cg \approx \Wbal$ and both reach their nadir together at exhaustion. During a steady
supra-\CP{} effort the two bars track \Wbal{} and each other; they diverge only in the activation ramp and
in recovery.
\item \textbf{Maximal:} the tapered ceiling makes \Rp{} decay $\sim$geometrically to its nadir (a
$\sim$10~s all-out sprint leaves the fast reserve at $\approx$20--30\%; \cref{subsec:behaviour}), and
fast-reserve dominance emerges from the ceiling, not the weight. This residual is a fast-reserve-of-\Wp{}
level, motivated by---but not validated against---muscle-PCr biopsy values.
\item \textbf{Energy conserved when a cap binds.} Residual $\mathrm{unmet}$ is banked in a deficit $D$, so
$(\Rp + \Rg - D)$ drops by exactly $\Delta\,\Delta t$ per second. Work booked to $D$ never drains a
reserve, so a large $D$ leaves the two bars reading slightly full; $D$ is an accounting term for the rare
supra-flux case, usually signalling a stale \Pos.
\end{itemize}

The depletion split is not identifiable from power: total draw above \CP{} is all the pedals see, and any
$(\fp, \gpmax)$ summing to the same total is observationally equivalent. In principle the bi-exponential
recovery of \Wp{} constrains the split; in practice that fit is ill-conditioned (the fast component is
nearly invisible at standard reconstitution sampling), so \fp{} is assumed and weakly constrained, not
recovered.

\paragraph{Case 2 --- restoration ($\Delta \le 0$, at/below \CP).} Each reserve refills exponentially; the
fast reserve is gated by oxidative headroom, and the slow reserve and deficit recover at an
intensity-dependent rate:

\begin{lstlisting}[style=pseudo]
g <- g * exp(-dt/tau_off)                        # glycolytic de-activation (tau_off = tau_on default)

# FAST RESERVE: resynthesis needs aerobic ATP ABOVE the ride's own demand, so it is gated by the
# OXIDATIVE HEADROOM (CP - P): full at rest, near-arrested at CP. gate_p scales the AMOUNT (not exponent).
gate_p = max(0, (CP - P) / CP)
R_p   += gate_p * (C_p - R_p) * (1 - exp(-dt/tau_p))

# SLOW RESERVE and deficit recover whenever P < CP at the intensity-dependent W'bal rate
# tau_W'(CP - P) = 546*exp(-0.01*(CP - P)) + 316, amplitude re-anchored so the 20 W passive
# rate reproduces the reconstitution curve. Bounded, and fitted across recovery powers.
if P < CP:
    f    = gate20 * tau_W'(CP - 20) / tau_W'(CP - P)   # = gate20 at 20 W anchor, < 1 above
    b    = f * (1 - exp(-dt/tau_g))
    R_g += (C_g - R_g) * b
    D   -= D * b
\end{lstlisting}

\emph{The intensity-dependent recovery shape.} A linear gate $(\LTone - P)/\LTone$ that falls to zero at
\LTone{} makes the effective recovery constant diverge as $P \to \LTone$; under such a gate the model
cannot complete a standard $4\times[\SI{4}{\minute}@\SI{330}{\watt} / \SI{4}{\minute}@\SI{150}{\watt}]$
session (see \cref{subsec:behaviour}). Adopting the established \Wp-recovery form
$\tauW(D_{\CP}) = 546\,e^{-0.01(\CP-P)}+316$\citep{Skiba2012}---already fitted across recovery
powers---for the shape, with its amplitude re-anchored so the 20~W passive rate reproduces the
reconstitution curve, yields bounded, monotone recovery that tracks the reference law within $\sim$15\%
everywhere and completes the session (\cref{tab:recovery}).

\emph{The oxidative-headroom gate.} $\gatep = (\CP-P)/\CP$ prevents the fast-reserve bar refilling while
the rider is under load. Its shape is a modelling choice: a linear, $\sqrt{\cdot}$, or sigmoid form swings
the 60~s post-effort recovery number by 21--31 points at \LTone{} (46\,/\,65\,/\,34\% after a
700~W/20~s effort). It also scales the refill only and adds no sub-\CP{} \emph{drain}, so a rider holding
0.95\,\CP{} for 30~min---never above \CP---keeps a full fast-reserve bar; the structurally correct
remedy is to relax \Rp{} toward an intensity-dependent steady state $\Rp^{*}(P)$, which is stated as a
limitation (\cref{sec:assumptions}). After a 20~s/700~W effort (fast reserve $\to$ 23\%), 60~s of recovery
restores it to 87\% at 0~W, 46\% at 0.8\,\CP, and 25\% at 0.99\,\CP.

The $g$ de-activation lets the activation ramp re-fire on each interval of a repeated-bout set (otherwise
$g$ ratchets to 1 on the first surge). $\tauoff = \tauon$ is assumed: activation, not de-activation, was
measured.\citep{Parolin1999} \LTone{} sets the anchor band of the recovery shape rather than a hard on/off
wall; recovery proceeds for any $P < \CP$, fastest at low intensity and slowing smoothly through the tempo
band. It is a measured threshold in watts; where only \CP{} is known, $\LTone \approx 0.80\,\CP$ is a
fallback and the weakest default in the model.

\paragraph{Optional realism (off by default).}
\begin{itemize}[leftmargin=1.4em,itemsep=1pt]
\item \emph{pH-slowed PCr recovery:} $\tau_{\mathrm{p,eff}} = \taup\,(1 + k(1 - \Rg/\Cg))$. It cannot be
pushed to reproduce muscle-PCr recovery ($\sim$78.7\% at 3.8~min\citep{Bogdanis1996}) without breaking the
reconstitution curve and the $4\times4$ (the $k \approx 16$ that hits the biopsy value drops the curve to
30/68/94 and the $4\times4$ to $-30$\%)---the same point as \cref{subsec:recon}: \Wp{} and muscle PCr do
not share a time constant.
\item \emph{Aerobic ramp} ($\tauaer \approx 25~\mathrm{s}$): a first-order aerobic supply $A(t)$ rising
toward $\min(P,\CP)$, with $\mathrm{need} = (P - A)\,\Delta t$, reproducing the onset O\textsubscript{2}
deficit; it also absorbs the shortfall across repeated sprints (\cref{subsec:repeatedsprint}).
\item \emph{Glycolytic flux fatigue} ($g_{\mathrm{fat}}$): scale the glycolytic ceiling down as the slow
reserve empties, $\mathrm{rate}_{g,\mathrm{eff}} = \mathrm{rate}_g\,(\Rg/\Cg)^{g_{\mathrm{fat}}}$,
capturing acidotic inhibition of phosphorylase/PFK (\cref{subsec:repeatedsprint}).
\end{itemize}

\subsection{Reported quantities}\label{subsec:reported}

\begin{itemize}[leftmargin=1.4em,itemsep=1pt]
\item \textbf{Fast-reserve level} $= \Rp/\Cp$ (0--100\%) --- ``punch'' remaining.
\item \textbf{Slow-reserve level} $= \Rg/\Cg$ (0--100\%) --- ``sustained dig'' remaining.
\item \textbf{Consumption rate} (per system, W) --- $\mathrm{take}_p/\Delta t$, $\mathrm{take}_g/\Delta t$.
\item \textbf{Combined \Wbal} $= (\Rp + \Rg - D)/\Wp$ --- energy-conserving in depletion (matches a
single-tank \Wbal{} reference on above-\CP{} segments) and deliberately divergent in recovery.
\item \textbf{Two distinct flags:} \texttt{rate\_limited} (power beyond the reserves' flux caps,
$\mathrm{unmet}>0$, reserves not empty---usually a stale \Pos, triggering a ``recalibrate sprint power''
hint) and \texttt{exhausted} (both reserves at/near zero). These are different conditions: a fresh rider
with a mis-set \Pos{} can be rate-limited but not exhausted.
\end{itemize}

\subsection{Model behaviour}\label{subsec:behaviour}

A short, very hard surge draws mostly from \Rp{} and refills afterward at a rate set by recovery intensity.
A maximal $\sim$10~s sprint leaves $\Rp \approx 20$--30\% (not empty) from the emergent depletion constant
$\taudep = \Cp/\Ppmax$; because \taudep{} swings $\sim$4--13~s across riders, the 10~s residual ranges from
$\sim$10\% to $\sim$46\%. A sustained supra-\CP{} effort draws from both reserves, declining together to
their nadir; the drawdown is front-loaded, not linear (450~W: 68/42/18/1 at quarter-TTE intervals). During
such steady efforts both bars are affine in \Wbal, so their only non-redundant content is the ramp offset
and the post-effort recovery transient.

For a constant supra-\CP{} effort, $\Rp/\Cp$ is asymptotically a function of the fraction $\varphi$ of \Wp{}
spent, but only for efforts long relative to \tauon{} and below the ceiling regime; the activation ramp
adds a second intensity dependence, so ``hold different supra-\CP{} powers to exhaustion and compare
fast-reserve-at-$\varphi$'' is a \emph{falsification} test (a pass means the bar is a deterministic
restatement of \Wbal{} in depletion), not a validation.

\Cref{tab:battery} reports a synthetic battery at $\CP = 255~\mathrm{W}$, $\Wp = 20~\mathrm{kJ}$,
$\fp = 0.25$, $\Ppmax = 690~\mathrm{W}$ ($\grate = 0.5$), $\taup = 27$, $\taug = 470$,
$\tauon = \tauoff = 6$, $\LTone = 204~\mathrm{W}$. \Cref{tab:recovery} shows the recovery law's behaviour
across recovery powers: the intensity-dependent shape is bounded and tracks the reference \Wp-recovery law
within $\sim$15\%, and the $4\times[\SI{4}{\minute}@330 / \SI{4}{\minute}@150]$ session completes
(\Wbal{} at rep starts $100/49/7/-24$\%).

\begin{table}[htbp]\centering\small
\caption{Synthetic verification battery.}
\label{tab:battery}
\renewcommand{\arraystretch}{1.2}
\begin{tabularx}{\linewidth}{@{}Xl@{}}
\toprule
\textbf{Test} & \textbf{Result} \\
\midrule
Sustained 450~W --- fast reserve at 25/50/75/100\% of TTE & 68 / 42 / 18 / $\sim$1\% (front-loaded) \\
Both reserves at exhaustion & fast $\sim$1\% / slow $\sim$0\% (empty together) \\
945~W ($=\Pos$) 10~s sprint --- fast-reserve residual & 23\% (ceiling-dominated) \\
$6\times[\SI{10}{\second}@700 / \SI{30}{\second}@150]$ --- fast reserve per bout & 49\,$\to$\,17\,$\to$\,10\,$\to$\,8\,$\to$\,8\,$\to$\,7\% \\
1200~W hold (caps bind) --- \Wbal{} conservation & leak $= 0$~J \\
Fast-reserve recovery after 700~W/20~s, 60~s @ \{0,128,204,242,252\}~W & 87 / 68 / 46 / 30 / 25\% \\
Debt repayment at 0.9\,\CP & $D$ clears slowly (recovery-law rate), not frozen \\
\bottomrule
\end{tabularx}
\end{table}

\begin{table}[htbp]\centering\small
\caption{\Wbal{} recovery half-time after full depletion, holding a constant recovery power: the
intensity-dependent shape tracks the reference law and stays bounded.}
\label{tab:recovery}
\renewcommand{\arraystretch}{1.2}
\begin{tabular}{@{}lccc@{}}
\toprule
\textbf{Recovery power} & \textbf{Linear-gate half-time} & \textbf{Intensity-dependent shape} & \textbf{Reference \tauW} \\
\midrule
20~W (anchor)      & 201~s & \textbf{201~s} & 255~s \\
100~W (39\%\,\CP)  & 353~s & \textbf{241~s} & 299~s \\
150~W (59\%\,\CP)  & 675~s & \textbf{291~s} & 351~s \\
190~W (75\%\,\CP)  & 2601~s & \textbf{367~s} & 417~s \\
204~W ($=\LTone$)  & $\infty$ & \textbf{409~s} & 446~s \\
\bottomrule
\end{tabular}
\end{table}

In depletion the model is a faithful generalisation of \Wbal{} ($(\Rp+\Rg-D)$ falls by exactly
$\Delta\,\Delta t$ per second); in recovery it is deliberately different (bi-exponential, intensity-gated),
so it does not reduce to a single \tauW. It is a generalisation in structure and depletion behaviour, not a
strict superset.

% =====================================================================
\section{Real-time on-device implementation}\label{sec:impl}

The per-second cost is a handful of multiplies and one $\exp()$ per reserve---trivial for a head unit. The
target design is a Garmin Connect~IQ full-screen \emph{Data Field} (\texttt{Toybox.WatchUi.DataField}),
called once per second (a natural fit for $\Delta t = 1~\mathrm{s}$), reading
\texttt{Activity.Info.currentPower}. Configuration (\CP, \Wp, \fp, \Ppmax, \LTone, and the recovery
constants) is exposed via Connect~IQ application properties. The display is two vertical bar gauges
(fast and slow reserve) with numeric \% and colour bands; per-system live consumption is a
\emph{modelled share}, not a reading, because when the ceiling is slack it is merely $(P - \CP)$ split by a
fixed fraction (\cref{sec:assumptions}), and is labelled as such or omitted. Both reserves and consumption
are written to the FIT file via \texttt{FitContributor} fields for post-ride analysis. The state is a
handful of scalars ($\Rp$, $\Rg$, the activation $g$, the deficit $D$)---no arrays---well within Connect~IQ
memory budgets.

% =====================================================================
\section{Validation strategy}\label{sec:validation}

The model is a hypothesis and must be checked against data. Its headline feature---a separate fast reserve
tracked in real cycling---is the hardest thing to validate, so we separate tests that check the plumbing
from the one test that would earn the second tank.

\begin{enumerate}[leftmargin=1.6em,itemsep=3pt,label=\textbf{\arabic*.}]
\item \textbf{Face validity.}\label{item:face} Synthetic power traces (single sprint, repeated sprints,
supra-\CP{} hold, sub-\CP{} recovery) should produce the qualitative behaviours of \cref{subsec:behaviour}.
\item \textbf{Backward compatibility (depletion only).}\label{item:backcompat} The combined balance
$(\Rp+\Rg-D)$ reproduces a reference \Wbal{} implementation in depletion (it drops by exactly
$\Delta\,\Delta t$ per second, deficit-preserved even when a cap binds). It will not match in recovery, by
construction, so the test is scoped to above-\CP{} segments.
\item \textbf{Reconstitution targets.}\label{item:recon} Combined recovery after a full \Wp{} depletion is
compared to the published curve (37/65/86\% at 2/6/15~min\citep{Ferguson2010}). The curve is blind to the
fast reserve (a $\taup \approx 27~\mathrm{s}$ process is $\sim$99\% complete by the first 2~min sample), so
it validates slow-reserve recovery and the size of the fast offset, not fast kinetics; and it fixes \taug{}
while constraining a $(\fp,\taug)$ ridge, not \fp{} alone (\cref{subsec:recon}).
\item \textbf{Fast-recovery kinetics.}\label{item:fastrecov} The load-bearing test for \taup{} is a
repeated-bout protocol with early recovery sampling (all-out efforts separated by 10/20/30/45/60~s), with
between-bout power recovery correlated with modelled fast-reserve recovery.\citep{Bogdanis1996} Standard
reconstitution data cannot supply this; it must be measured deliberately.
\item \textbf{System-specific truth, and its limit.}\label{item:31p} Blood-lactate kinetics can corroborate
the slow reserve, but the fast reserve's gold standard, \degC, is essentially unavailable in real cycling
(the muscle must be inside a magnet), so it exists only for small-muscle knee-extension or immediate
post-exercise. The headline novelty cannot be checked against \degC{} in the modality it targets.
\item \textbf{The head-to-head that earns the second tank.}\label{item:headtohead} Because the model is a
near-superset of single-tank in depletion, ``it fits better'' is guaranteed and is not evidence the
compartments are real. The decisive test is a head-to-head against single-tank \Wbal{} on an out-of-sample
intermittent-tolerance prediction.\citep{Weigend2021} Since depletion here is single-tank by construction,
the only lever is the recovery law, so this is a test of the recovery law, not of the two-compartment
hypothesis; the protocol should choose recovery-valley powers straddling \LTone, where the two recovery
laws diverge most.
\item \textbf{Field calibration.}\label{item:field} Fit \Ppmax{} and the $\tau$'s per athlete to maximal
and repeated-bout tests. \fp{} is not a routine target (the split is unidentifiable from power); any
single-tank comparator in the head-to-head must have its own \tauW{} fit to the same data, or the dual-tank
wins trivially.
\end{enumerate}

\subsection{Empirical usefulness on real interval rides}\label{subsec:usefulness}

Two data-in-hand questions decide whether the feature is worth building. Both were evaluated on six real
interval / VO\textsubscript{2}max sessions (the model's best case; a steady endurance ride would score far
lower).

\emph{Is the fast-reserve bar informative, and is it a function of \Wbal?} The bar is below 95\% (live)
for a per-ride 20 / 31 / 38 / 52 / 61 / 79\% of ride time (median $\sim$45\%). Binning \Wbal{} and
measuring the spread of $\Rp/\Cp$ within each 1-point bin gives a mean spread of $\sim$58 points (max
$\sim$93); regressing the bar on its best affine \Wbal{} predictor gives $R^2 \approx 0.03$, with the bar
$>$5 points off that affine line $\sim$93\% of the time. The bar is emphatically not a function of \Wbal.

\emph{Does the depletion machinery change the display, versus a recovery-only null?} Running the full
model and the null on the same traces, the fast-reserve bar diverges by up to 25--60 points, with maxima at
the starts of recovery valleys. Divergence is not accuracy, however: two models differing by 25--60 points
tells us the caps \emph{change} the display, not that the changed display is \emph{right}. These results
establish that the bar carries information; whether that information is correct is
\cref{item:headtohead}\,/\,\cref{subsec:repeatedsprint}.

\subsection{Training-load partitioning}\label{subsec:trainingload}

The cumulative per-system load over a session---how many kJ came from the fast versus the slow
system---is a different functional of the power trace than \Wbal, and is the strongest reason to build the
second tank. For two fully-specified archetypes (alactic:
$10\times[\SI{6}{\second}@900 / \SI{300}{\second}@150]$; glycolytic:
$5\times[\SI{60}{\second}@360 / \SI{120}{\second}@150]$), reporting the deficit $D$ as a third quantity
(\cref{tab:partition}), the model separates the sessions by $\sim$39 points. Three qualifications are
essential.

\begin{table}[htbp]\centering\small
\caption{Per-system load for two session archetypes (default parameters), with the deficit reported.}
\label{tab:partition}
\renewcommand{\arraystretch}{1.2}
\begin{tabular}{@{}lcccc@{}}
\toprule
\textbf{Session} & \textbf{fast (kJ)} & \textbf{slow (kJ)} & \textbf{$D$ (kJ)} & \textbf{fast \% of (fast+slow)} \\
\midrule
$10\times[\SI{6}{\second}@900]$ --- alactic     & 29.2 & 8.7  & 1.2 & \textbf{77\%} \\
$5\times[\SI{60}{\second}@360]$ --- glycolytic   & 13.2 & 21.5 & 1.8 & \textbf{38\%} \\
\bottomrule
\end{tabular}
\end{table}

\textbf{(a) The mechanism is the glycolytic flux ceiling scaled by the ramp}, not the ramp alone. Ablating
the model on the alactic session (\cref{tab:ablation}): removing the ceilings costs $\sim$30 points,
removing the ramp $\sim$25; neither alone reaches the full 77\%. It is an interaction---during a 6~s sprint
the glycolytic ceiling caps glycolysis at a few tens of watts while $\sim$645~W is demanded, and the spill
goes to the fast reserve---so the load-bearing parameter is \grate. Correspondingly, a 20$\times$ change in
\tauoff{} moves the split by only 1--2 points, whereas \grate{} $0.3\rightarrow0.7$ moves the separation
48$\rightarrow$33 points.

\begin{table}[htbp]\centering\small
\caption{Ablation on the alactic session: the ceiling and the ramp both contribute; neither alone reaches
the full separation.}
\label{tab:ablation}
\renewcommand{\arraystretch}{1.2}
\begin{tabular}{@{}lcc@{}}
\toprule
\textbf{Configuration} & \textbf{alactic fast \%} & \textbf{separation vs glycolytic} \\
\midrule
full model                     & \textbf{77\%} & $\sim$39 pts \\
no activation ramp ($g\equiv1$) & 52\%          & $\sim$15 pts \\
no rate ceilings (share rule only) & 47\%       & $\sim$9 pts \\
neither                        & 25\%          & $-12$ pts \\
\bottomrule
\end{tabular}
\end{table}

\textbf{(b) The metric detects supra-flux-ceiling power, so it is blind to submaximal alactic work.}
Holding the $10\times6$~s structure fixed and sweeping power, the reading climbs
52\,$\to$\,62\,$\to$\,66\,$\to$\,72\,$\to$\,77\% at 450/550/600/700/900~W. A textbook alactic session at
450~W reads $\sim$52\%, close to the glycolytic session's 38\%.

\textbf{(c) The honest baseline is a stopwatch, not \Wbal.} A coach with the power file and no model can
compute the fraction of supra-\CP{} work done in efforts under 15~s---100\% versus 0\% for these two
sessions, perfect separation with zero parameters. The open question is therefore whether the per-system
load carries information beyond effort duration and intensity (regressing fast-reserve \% on a
duration/intensity baseline across real rides answers it). What does survive: the separation is robust to
every assumption it rests on (including substituting Ferguson's $\taup = \SI{107}{\second}$, $<$1~pt), and
it is a genuinely different functional of the power trace. Absolute kJ scale with \fp, $D$ must be reported
(and the split refused when $D$ is a large fraction of total---e.g.\ a $4\times4$ books $\sim$45\% of its
anaerobic work to $D$), and this is a training-design signal, not a validated training outcome.

\subsection{A compartment-level test the model fails}\label{subsec:repeatedsprint}

Muscle biopsy during cycle ergometry \emph{is} in-modality, compartment-level data, and repeated-sprint
studies partition ATP supply by system from vastus-lateralis biopsies during maximal cycle-ergometer
sprints.\citep{Gaitanos1993,SciRep2024} The model's alactic protocol ($10\times6$~s maximal sprints) is
exactly this design, and the biopsy literature reports the glycolytic ATP share falling from $\approx$40\%
in the first sprint to $<$10\% in the last, with aerobic metabolism rising to cover the shortfall. The
model fails this test in level and direction (\cref{tab:repeatsprint}).

\begin{table}[htbp]\centering\small
\caption{Repeated-sprint glycolytic share: the model, as formulated, runs at the wrong level and in the
wrong direction relative to biopsy data.}
\label{tab:repeatsprint}
\renewcommand{\arraystretch}{1.2}
\begin{tabular}{@{}lccc@{}}
\toprule
 & \textbf{Sprint 1} & & \textbf{Sprint 10} \\
\midrule
Observed\citep{Gaitanos1993,SciRep2024} --- glycolytic share & $\sim$40\% & $\searrow$ & $<$10\% \\
Model, default ($g_{\mathrm{fat}}=0$) & 23\% & $\nearrow$ & 46\% \\
Model, $g_{\mathrm{fat}}=1$ + aerobic ramp & 23\% & $\rightarrow$ & $\sim$34\% (flat, then falling) \\
\bottomrule
\end{tabular}
\end{table}

The model starts glycolysis at about half the observed share and has it \emph{rise} across bouts, because
it has the two fatigue mechanisms backwards: the fast-reserve ceiling ($\mathrm{rate}_p = \Ppmax\,\Rp/\Cp$)
fades as the reserve empties, so more demand spills to glycolysis bout over bout, while glycolytic flux has
no fatigue term. The optional $g_{\mathrm{fat}}$ term plus the aerobic ramp flatten the trend and pull the
late bouts down but do not reproduce 40\%$\rightarrow$$<$10\%; reaching 40\% in sprint~1 would require a
higher \grate{} that then wrecks the separation of \cref{subsec:trainingload}. Two conclusions follow. This
is a genuinely falsifiable result: a number wrong in a specific, fixable way, obtained because the
partition finally faced compartment-level data. And it fixes the status of the output: the per-system split
is a \Wp-decomposition heuristic and a descriptive statistic of a power file, not a validated bioenergetic
ATP partition. Calibrating \grate{} to the sprint-1 biopsy value and fitting $g_{\mathrm{fat}}$ to the
decay---using data already in the literature---is the highest-value experiment remaining and the one test
that could directly falsify or earn the two-compartment hypothesis.

% =====================================================================
\section{Assumptions and limitations}\label{sec:assumptions}

\begin{itemize}[leftmargin=1.4em,itemsep=3pt]
\item \textbf{The split \fp{} is assumed and only consistent with (not measured by) the recovery data.}
Power cannot identify the depletion split; the reconstitution curve constrains a $(\fp,\taug)$ ridge, and
physiology selects \fp{} on it (\cref{subsec:recon}). Supported band $\sim$0.20--0.25; outputs are
sensitive to it.
\item \textbf{The reserves are a decomposition of \Wp, not latent metabolic state.} When the fast reserve
is full (rest), the slow bar is an affine transform of \Wbal. Under the recovery gate the effective
fast-reserve constant is $\taup/\gatep$---27~s only at a standstill, rising to $\sim$44~s at 100~W,
$\sim$66~s at 150~W, $\sim$135~s at \LTone---so while riding the bar is live for a substantial fraction of
interval-ride time; the ``affine'' caveat applies to stopped or easy recovery, not to riding.
\item \textbf{\gatep{} is an assumption with an undefended shape} (linear vs $\sqrt{\cdot}$ vs sigmoid
swings the post-effort recovery number 21--31 points at \LTone) and adds no sub-\CP{} drain.
\item \textbf{\grate{} is load-bearing} (it generates the training-load separation) and has never been fit
to sprint biopsy data (\cref{subsec:repeatedsprint}).
\item \textbf{The fast-reserve depletion rate $\taudep = \Cp/\Ppmax$ is assumed}, is the ratio of three
parameters, and swings $\sim$3$\times$ across plausible riders ($\sim$4--13~s); it should be sanity-checked
against literature PCr depletion half-times.
\item \textbf{\tauoff{} is assumed and not load-bearing} (a 20$\times$ change moves the training-load split
1--2 points). The discriminating content is the glycolytic flux ceiling scaled by the ramp, not \tauoff.
\item \textbf{\LTone{} should be measured, not derived from \CP.} The $0.80\,\CP$ fallback is the weakest
default and materially changes recovery estimates.
\item \textbf{Small modelling choices are assumptions, not derivations:} the recovery shape (an
established prior fitted across recovery powers, but still a choice of form); lumping deficit clearance and
slow-reserve refill onto the same \taug; and the spill order (tested inert on the sessions that matter).
\item \textbf{The deficit $D$ corrupts the display at maximal effort} (large $D$ leaves the bars reading
slightly full). Its late-effort recurrence is driven by the fullness taper collapsing the flux ceiling
below what a deep-but-still-sprinting rider produces---structural, not a settings error; $D$ should be
reported and the per-system split refused when $D$ is a large fraction of total.
\item \textbf{Identifiability degeneracies must be named before fitting:} \fp{} and \Ppmax{} both set the
sprint residual; \taup{} and the pH coefficient both set the repeated-bout decay; \grate{} and the ramp
both set the partition level.
\item \textbf{Per-system live consumption is a modelled share, not a reading:} where the ceiling is slack,
$\mathrm{take}_p = \Delta\,\Cp/(\Cp+\Cg g)$ contains neither reserve and settles to a fixed $\Delta\,\fp$
split---$(P-\CP)$ rescaled.
\item \textbf{A structural product fork.} The lightweight-model rationale serves the \emph{pacing} use
case on a head unit. The training-load metric (\cref{subsec:trainingload}) is a post-hoc computation on a
power file; off-device, a better-resolved model (hydraulic or bioenergetic) becomes available and can be
calibrated to biopsy data. A live two-bar pacing display and a post-ride training metric have different
centres of gravity.
\item \textbf{Fixed time constants} ignore the documented pH-dependence and bout-to-bout slowing of
fast-reserve recovery unless the optional term is enabled.
\item \textbf{A hard \CP{} boundary} omits the aerobic slow component and onset kinetics unless the aerobic
ramp is enabled; expect over-attribution to the anaerobic reserves in long low-intensity recovery.
\item \textbf{Power is whole-body and mechanical}, not muscle-specific; the model cannot see local fatigue,
cadence, or fibre-type effects. It is a training/pacing aid, \textbf{not medical or diagnostic}, and
parameters drift with fitness, fatigue, heat, and altitude.
\end{itemize}

% =====================================================================
\section{Conclusion}\label{sec:conclusion}

The single-tank \Wp{} model earned its place by being power-native and light enough to run anywhere, but it
answers ``how much anaerobic work is left?'' without answering ``in which system?'' The dual-tank model
offers a physiologically-motivated split at negligible computational cost: divide \Wp{} into a fast,
rate-limited, fast-recovering reserve and a slow, large, slowly-recovering reserve; share supra-\CP{}
demand by capacity; cap each by its peak flux (the fast reserve tapering with fullness); and refill them
with system-specific, intensity-gated laws. It generalises \Wbal{} in depletion, diverges from it by design
in recovery, runs at 1~Hz on a head unit, and surfaces a two-system decomposition---\emph{punch left} and
\emph{dig left}---most informative in the transients where pacing decisions concentrate.

Two framing points are essential to state plainly. The reserves are compartments of \Wp, not of muscle
metabolites: \Wp{} recovers $\sim$4$\times$ faster than the PCr and lactate the tanks are named after
(\cref{subsec:recon}), which makes the time constants correct for \Wp{} while the metabolic labels remain
motivation. And the per-system split is falsifiable in-modality: run against repeated-sprint biopsy data it
fails in level and direction (\cref{subsec:repeatedsprint}), so it is a descriptive \Wp-statistic, not a
validated bioenergetic partition. The model is therefore best located as a lightweight, honest heuristic
decomposition of \Wp---with a live pacing display whose value rests on the fast-reserve recovery transient,
and a post-ride training metric that is promising, descriptive, and, for the first time, has a data-in-hand
test that can earn or sink it. The remaining moves are experiments, not arguments: the recovery-law
head-to-head (\cref{item:headtohead}) and the biopsy calibration (\cref{subsec:repeatedsprint}).

% =====================================================================
\begin{thebibliography}{99}
\setlength{\itemsep}{2pt}

\bibitem{Skiba2012} Skiba PF, Chidnok W, Vanhatalo A, Jones AM. Modeling the expenditure and reconstitution
of work capacity above critical power. \emph{Med Sci Sports Exerc}. 2012;44(8):1526--1532.
\href{https://doi.org/10.1249/MSS.0b013e3182517a80}{doi:10.1249/MSS.0b013e3182517a80}.
\href{https://pubmed.ncbi.nlm.nih.gov/22382171/}{PMID:22382171}.

\bibitem{Margaria1933} Margaria R, Edwards HT, Dill DB. The possible mechanisms of contracting and paying
the oxygen debt and the r\^ole of lactic acid in muscular contraction. \emph{Am J Physiol}.
1933;106(3):689--715. \href{https://doi.org/10.1152/ajplegacy.1933.106.3.689}{doi:10.1152/ajplegacy.1933.106.3.689}.

\bibitem{diPrampero1999} di Prampero PE, Ferretti G. The energetics of anaerobic muscular metabolism: a
reappraisal of older and recent concepts. \emph{Respir Physiol}. 1999;118(2--3):103--115.
\href{https://doi.org/10.1016/S0034-5687(99)00083-3}{doi:10.1016/S0034-5687(99)00083-3}.
\href{https://pubmed.ncbi.nlm.nih.gov/10647856/}{PMID:10647856}.

\bibitem{Ferguson2010} Ferguson C, Rossiter HB, Whipp BJ, Cathcart AJ, Murgatroyd SR, Ward SA. Effect of
recovery duration from prior exhaustive exercise on the parameters of the power-duration relationship.
\emph{J Appl Physiol}. 2010;108(4):866--874.
\href{https://doi.org/10.1152/japplphysiol.91425.2008}{doi:10.1152/japplphysiol.91425.2008}.
\href{https://pubmed.ncbi.nlm.nih.gov/20093659/}{PMID:20093659}.

\bibitem{Harris1976} Harris RC, Edwards RHT, Hultman E, Nordesj\"o LO, Nylind B, Sahlin K. The time course
of phosphorylcreatine resynthesis during recovery of the quadriceps muscle in man. \emph{Pfl\"ugers Arch}.
1976;367(2):137--142. \href{https://doi.org/10.1007/BF00585149}{doi:10.1007/BF00585149}.
\href{https://pubmed.ncbi.nlm.nih.gov/1034909/}{PMID:1034909}.

\bibitem{Yoshida2013} Yoshida T, Watari H, Tagawa K. Effects of active and passive recoveries on muscle
metabolism during exercise as studied by \textsuperscript{31}P-MRS. \emph{Scand J Med Sci Sports}.
2013;23(5):e320--e330. \href{https://doi.org/10.1111/sms.12081}{doi:10.1111/sms.12081}.
\href{https://pubmed.ncbi.nlm.nih.gov/23662804/}{PMID:23662804}.

\bibitem{Gaitanos1993} Gaitanos GC, Williams C, Boobis LH, Brooks S. Human muscle metabolism during
intermittent maximal exercise. \emph{J Appl Physiol}. 1993;75(2):712--719.
\href{https://doi.org/10.1152/jappl.1993.75.2.712}{doi:10.1152/jappl.1993.75.2.712}.
\href{https://pubmed.ncbi.nlm.nih.gov/8226455/}{PMID:8226455}.

\bibitem{SciRep2024} Repeated-sprint ATP partitioning by energy system from vastus-lateralis biopsy during
cycle ergometry. \emph{Sci Rep}. 2024;14.
\href{https://doi.org/10.1038/s41598-024-78916-z}{doi:10.1038/s41598-024-78916-z}.

\bibitem{Bogdanis1996} Bogdanis GC, Nevill ME, Boobis LH, Lakomy HKA. Contribution of phosphocreatine and
aerobic metabolism to energy supply during repeated sprint exercise. \emph{J Appl Physiol}.
1996;80(3):876--884. \href{https://doi.org/10.1152/jappl.1996.80.3.876}{doi:10.1152/jappl.1996.80.3.876}.
\href{https://pubmed.ncbi.nlm.nih.gov/8964751/}{PMID:8964751}.

\bibitem{Parolin1999} Parolin ML, Chesley A, Matsos MP, Spriet LL, Jones NL, Heigenhauser GJF. Regulation
of skeletal muscle glycogen phosphorylase and PDH during maximal intermittent exercise. \emph{Am J Physiol
Endocrinol Metab}. 1999;277(5):E890--E900.
\href{https://doi.org/10.1152/ajpendo.1999.277.5.E890}{doi:10.1152/ajpendo.1999.277.5.E890}.
\href{https://pubmed.ncbi.nlm.nih.gov/10567017/}{PMID:10567017}.

\bibitem{GonzalezAlonso2000} Gonz\'alez-Alonso J, Quistorff B, Krustrup P, Bangsbo J, Saltin B. Heat
production in human skeletal muscle at the onset of intense dynamic exercise. \emph{J Physiol}.
2000;524(2):603--615. \href{https://doi.org/10.1111/j.1469-7793.2000.00603.x}{doi:10.1111/j.1469-7793.2000.00603.x}.
\href{https://pubmed.ncbi.nlm.nih.gov/10766936/}{PMID:10766936}.

\bibitem{Skiba2015} Skiba PF, Fulford J, Clarke DC, Vanhatalo A, Jones AM. Intramuscular determinants of
the ability to recover work capacity above critical power. \emph{Eur J Appl Physiol}.
2015;115(4):703--713. \href{https://doi.org/10.1007/s00421-014-3050-3}{doi:10.1007/s00421-014-3050-3}.
\href{https://pubmed.ncbi.nlm.nih.gov/25425383/}{PMID:25425383}.

\bibitem{Bartram2018} Bartram JC, Thewlis D, Martin DT, Norton KI. Accuracy of \Wp{} recovery kinetics in
high-performance cyclists---modeling intermittent work capacity. \emph{Int J Sports Physiol Perform}.
2018;13(6):724--728. \href{https://doi.org/10.1123/ijspp.2017-0034}{doi:10.1123/ijspp.2017-0034}.
\href{https://pubmed.ncbi.nlm.nih.gov/29140148/}{PMID:29140148}.

\bibitem{Morton1986} Morton RH. A three component model of human bioenergetics. \emph{J Math Biol}.
1986;24(4):451--466. \href{https://doi.org/10.1007/BF01236892}{doi:10.1007/BF01236892}.
\href{https://pubmed.ncbi.nlm.nih.gov/3805974/}{PMID:3805974}.

\bibitem{Weigend2021} Weigend FC, Siegler J, Brueckner D. A hydraulic model outperforms work-balance models
for predicting recovery kinetics from intermittent exercise. \emph{arXiv} preprint. 2021.
\href{https://arxiv.org/abs/2108.04510}{arXiv:2108.04510}.

\bibitem{Chorzempa2023} A dynamic bioenergetic model for intermittent cycling. \emph{Eur J Appl Physiol}.
2023;123. \href{https://doi.org/10.1007/s00421-023-05256-7}{doi:10.1007/s00421-023-05256-7}.
\href{https://pubmed.ncbi.nlm.nih.gov/37369795/}{PMID:37369795}.

\bibitem{Bangsbo1990} Bangsbo J, Gollnick PD, Graham TE, Juel C, Kiens B, Mizuno M, Saltin B. Anaerobic
energy production and O\textsubscript{2} deficit-debt relationship during exhaustive exercise in humans.
\emph{J Physiol}. 1990;422:539--559.
\href{https://doi.org/10.1113/jphysiol.1990.sp018000}{doi:10.1113/jphysiol.1990.sp018000}.
\href{https://pubmed.ncbi.nlm.nih.gov/2352192/}{PMID:2352192}.

\bibitem{ChorleyLamb2020} Chorley A, Lamb KL. The application of critical power, the work capacity above
critical power (\Wp), and its reconstitution: a narrative review. \emph{Sports (Basel)}. 2020;8(9):123.
\href{https://doi.org/10.3390/sports8090123}{doi:10.3390/sports8090123}.
\href{https://pubmed.ncbi.nlm.nih.gov/32899879/}{PMID:32899879}.

\bibitem{Bogdanis1998} Bogdanis GC, Nevill ME, Lakomy HKA, Boobis LH. Power output and muscle metabolism
during and following recovery from 10 and 20~s of maximal sprint exercise in humans. \emph{Acta Physiol
Scand}. 1998;163(3):261--272. \href{https://doi.org/10.1046/j.1365-201x.1998.00378.x}{doi:10.1046/j.1365-201x.1998.00378.x}.
\href{https://pubmed.ncbi.nlm.nih.gov/9715738/}{PMID:9715738}.

\end{thebibliography}

\end{document}
